\documentclass[10pt]{article}
\usepackage[T1]{fontenc}
\usepackage[utf8]{inputenc}
\usepackage[french]{babel}
\usepackage{amsmath,amssymb,mathrsfs,mathpazo}
\usepackage{geometry}
\geometry{a4paper,margin=2cm}
\usepackage{graphicx,array,tabularx,tabularray,colortbl,multirow,multicol,booktabs,enumitem,pifont,ulem,fancybox,diagbox,cancel}
\usepackage{pstricks,pstricks-add,pst-plot,pst-tree,pst-node,pst-text,pst-func,pst-math,pst-bspline,pst-3dplot,pst-eucl,pst-solides3d}
\let\pscalculate\undefined
\usepackage{pst-bezier}
\usepackage{tikz,pgfplots,tkz-tab}
\usetikzlibrary{arrows,arrows.meta,patterns,calc,positioning,trees,decorations.markings,shapes.geometric,shapes.misc,shapes.arrows,intersections,decorations.pathmorphing,decorations.pathreplacing,backgrounds,shadows,babel}
\pgfplotsset{compat=1.18}
\usepgfplotslibrary{fillbetween}
\usepackage[np]{numprint}
\usepackage{listings,makecell,pas-tableur,twoopt}
\let\barre\undefined
\usepackage[pstricks]{tablvar}
\definecolor{darkgreen}{RGB}{0,100,0}
\providecommand{\red}{\color{red}}
\providecommand{\blue}{\color{blue}}
\providecommand{\green}{\color{green}}
\providecommand{\black}{\color{black}}
\providecommand{\white}{\color{white}}
\providecommand{\gray}{\color{gray}}
\providecommand{\cyan}{\color{cyan}}
\providecommand{\magenta}{\color{magenta}}
\providecommand{\yellow}{\color{yellow}}
\providecommand{\violet}{\color{violet}}
\providecommand{\orange}{\color{orange}}
\providecommand{\euro}{\texteuro}
\providecommand{\EUR}{\texteuro}
\providecommand{\ieme}{\textsuperscript{e}}
\providecommand{\degres}{\ensuremath{^\circ}}
\providecommand{\degree}{\ensuremath{^\circ}}
\providecommand{\vv}[1]{\overrightarrow{#1}}
\providecommand{\Coord}[1]{#1}
\providecommand{\decosix}{}
\providecommand{\starredbullet}{\ensuremath{\bullet}}
\pagestyle{empty}
\setlength{\parindent}{0pt}
\usepackage{adjustbox}\DeclareUnicodeCharacter{2236}{:}\DeclareUnicodeCharacter{2190}{\ensuremath{\leftarrow}}
\providecommand{\R}{\mathbb{R}}
\providecommand{\N}{\mathbb{N}}
\providecommand{\Z}{\mathbb{Z}}
\providecommand{\Q}{\mathbb{Q}}
\providecommand{\D}{\mathbb{D}}
\providecommand{\C}{\mathbb{C}}
\providecommand{\e}{{\mathrm{e}}}
\providecommand{\vect}[1]{\overrightarrow{#1}}
\providecommand{\vectt}[1]{\overrightarrow{\mathrm{#1}}}
\providecommand{\vvt}[1]{\overrightarrow{\mathrm{#1}}}
\providecommand{\np}[2][]{{#2~#1}}
\providecommand{\vv}[1]{{\overrightarrow{#1}}}
\renewcommand{\d}{{\mathrm{d}}}
\providecommand{\degres}{{^\circ}}
\providecommand{\degre}{{^\circ}}
\providecommand{\blue}{\color{blue}}
\providecommand{\red}{\color{red}}
\providecommand{\enskip}{\quad}
\providecommand{\llbracket}{\lbrack\!\lbrack}
\providecommand{\rrbracket}{\rbrack\!\rbrack}
\providecommand{\euro}{{\text{€}}}
\providecommand{\no}{{\text{n°}}}
\providecommand{\interval}[2]{{[#1;#2]}}
\providecommand{\intervFO}[2]{{[#1;#2[}}
\providecommand{\dotfill}{\ldots}

\providecommand{\ds}{\displaystyle}
\providecommand{\textding}[1]{#1}
\providecommand{\Oij}{\(\mathrm{O};\vec{\imath},\vec{\jmath}\)}
\providecommand{\Oijk}{\(\mathrm{O};\vec{\imath},\vec{\jmath},\vec{k}\)}
\providecommand{\Ouv}{\(\mathrm{O};\vec{u},\vec{v}\)}
\providecommand{\intervFF}[2]{\ensuremath{[#1;#2]}}
\renewcommand{\euro}{\ifmmode\text{\texteuro}\else\texteuro\fi}
\DeclareUnicodeCharacter{00B0}{\ifmmode^{\circ}\else\textdegree\fi}
\DeclareUnicodeCharacter{20AC}{\ifmmode\text{\texteuro}\else\texteuro\fi}
% Énoncés : Éducation nationale. Transcription et corrigés : APMEP, Denis Vergès et contributeurs. Découpage et classement : CoopMaths, Sébastien Lozano, Éric Elter et contributeurs. Adaptation de la présentation : Maths Explorer.
% CC BY-SA : https://coopmaths.fr/www/about/
% https://coopmaths.fr/alea/static/bac/2023/tex/bac_2023_03_sujet2_etrangers_groupe1_2_cor.tex
\begin{document}
\noindent\textit{APMEP / CoopMaths. CC BY-SA. Présentation adaptée par Maths Explorer.}\par\medskip
\bigskip
\textbf{Partie A}

On considère la fonction $g$ définie sur $ ]- 1,5~;~+\infty[ $ par $g(x)=f(x)-x=\ln(2x+3)-1-x$.

\begin{enumerate}
\item $\lim\limits_{x\to-1,5^+}{2x+3}=0^+$ ( $x> - 1,5$ donc $2x+3>0$)

or $\lim\limits_{X\to 0}\ln (X) =-\infty$ donc, par composition, $\lim\limits_{x\to - 1,5}\ln(2x+3)= -\infty$.

De plus, $\lim\limits_{x\to -1,5}-1-x= 0,5$ donc, par somme,  $\lim\limits_{x\to -1,5}g(x)=-\infty$.

\item $g$ est dérivable sur $]-1,5~;~+\infty[$ comme composée de fonctions dérivables sur cet intervalle.

$g$ est de la forme $\ln(u)+v$ avec $u(x)=2x-3$ et $v(x)=-1-x$ .  On a donc $u'(x)=2$ et $v'(x)=-1$.

$g'=\dfrac{u'}{u}+v'$ donc, pour tout réel $x \in ]-1,5~;~+\infty[$, $g'(x) =\dfrac{2}{2x+3}-1=\dfrac{2-2x-3}{2x+3}=\dfrac{-2x-1}{2x+3}$.

Sur l'intervalle $]-1,5~;~+\infty[ $,  $2x+3$ est strictement positif donc le signe de $g'(x)$ est le signe de $-2x-1$ : $g'(x)$ est donc strictement positif sur l'intervalle  $ ]-1,5~;~-0,5[$ et strictement négatif sur l'intervalle  $ ]-0,5~;~+\infty[$.

$g$ est donc strictement croissante sur l'intervalle  $ ]- 1,5~;~-0,5[ $ et strictement décroissante sur l'intervalle  $ ]-0,5~;~+\infty[$.

\item
	\begin{enumerate}
		\item Sur l'intervalle $]-0,5~;~+\infty[ $, la fonction $g$ est une fonction continue (car dérivable) et strictement décroissante. D'où le tableau de variations :
		
\begin{center}
\psset{unit=1cm}
\begin{pspicture}(7,3.5)
\psframe(7,3.5)\psline(0,3)(7,3) \psline(1,0)(1,3.5)\psline(0,2.5)(7,2.5)
\uput[u](0.5,2.9){$x$} \uput[u](1.4,2.9){$-1,5$} \uput[u](4,2.9){$-0,5$} \uput[u](6.5,2.9){$+ \infty$}
\rput(0.5,2.75){$g'(x)$}\rput(0.5,1.25){$g$}
\uput[u](6.5,0){$- \infty$}\uput[d](4,2.5){$\approx 0,19$}\uput[u](1.5,0){$- \infty$}
\psline{->}(1.5,0.5)(3.5,2)\psline{->}(4.5,2)(6.5,0.5)
\rput(2.5,2.75){$+$}\rput(4,2.75){0}\rput(5.5,2.75){$-$}
\end{pspicture}
\end{center}
$g(-0,5) =\ln(2\times (-0,5)+3)- 1 + 0,5 = \ln(2) - 0,5 \approx 0,19$ et $\lim\limits_{x\to +\infty}g(x)=- \infty$ (admis), donc 0  est une valeur intermédiaire entre  $g(-0,5)$ et $\lim\limits_{x\to +\infty}g(x)$.
		
En vertu du corollaire au théorème des valeurs intermédiaires appliqué aux fonctions strictement monotones, l'équation $g(x)=0$ admet une unique solution sur l'intervalle $]-0,5~;~+\infty[$.
	\item La calculatrice donne :
	
$\bullet~$$g(0,2) \approx 0,24$ et $g(0,3) \approx - 0,02$ donc $0,2 < \alpha < 0,3$ ;
	
$\bullet~$$g(0,25) \approx 0,003$ et $g(0,26) \approx - 0,002$ donc $0,25\leqslant\alpha \leqslant 0,26$.
	\end{enumerate}
\end{enumerate}

\bigskip

\textbf{Partie B}

\medskip

\begin{enumerate}
\item $f$ est  croissante sur l'intervalle  $ ]-1,5~;~+\infty[ $  et donc en particulier sur l'intervalle $[-1~; \alpha]$ .

Donc $-1 \leqslant x \leqslant \alpha \Longrightarrow f(-1) \leqslant f(x) \leqslant f(\alpha)$.

Or $g(\alpha ) = 0 = f(\alpha) -\alpha$ donc $ f(\alpha)=\alpha$ et $f(-1) = \ln(2\times(-1)+3)-1= \ln(1) - 1 =- 1$

donc $- 1 \leqslant f(x) \leqslant\alpha $, c'est à dire $f(x) \in [-1~;~\alpha]$.
\item
	\begin{enumerate}
		\item
\emph{Initialisation :} à l'indice $n=0$ :
			
On a  $u_0=0$ et $u_1=f(u_0)=f(0)=\ln(3)-1\approx 0,099$  donc $-1 \leqslant u_0 \leqslant u_1 \leqslant\alpha$.

La propriété est vraie pour $n=0$.

\emph{Hérédité :} supposons que, pour un entier naturel $n$ donné, l'inégalité  $-1 \leqslant u_n \leqslant u_{n+1} \leqslant\alpha$ est vraie. Montrons que l'inégalité sera vraie au rang suivant :

Par hypothèse de récurrence : 

$\begin{aligned}[t]  & -1 \leqslant u_n \leqslant u_{n+1} \leqslant\alpha \\
	&\implies f(-1) \leqslant f(u_n) \leqslant f(u_{n+1}) \leqslant f(\alpha) \quad \text{ car } f \text{ est croissante sur  } [-1~;~\alpha]\\
	&\implies -1\leqslant u_{n+1}\leqslant u_{n+2} \leqslant\alpha  \quad \text{car } f(-1)=-1 \text{ et } f(\alpha)=\alpha\\ \end{aligned}$

Si l'inégalité est vérifiée à l'indice $n$,  alors, elle l'est aussi au rang suivant.

\emph{Conclusion :} L'inégalité est vérifiée à l'indice 0 et sa véracité est héréditaire pour tout indice $n$ naturel, donc, par principe de récurrence, on a  $-1 \leqslant u_n \leqslant u_{n+1} \leqslant\alpha$ pour tout entier naturel $n$.

	\item On vient d'établir que :

pour tout entier naturel $n$ , $u_n \leqslant\alpha$ donc la suite est majorée par $\alpha$

et  pour tout entier naturel $n$, $u_n \leqslant u_{n+1}$ donc la suite est  croissante, 

elle est donc convergente.
	\end{enumerate}
\end{enumerate}

\bigskip

\end{document}
